\section{El término constante}\label{sec:constante}

En la fórmula del predictor aparecen términos de tamaños muy distintos. Para identificar la
constante hay que conservar, en cada uno, todo lo que sobrevive después de restar $\log n!+n\Phi$.
Por ejemplo, la sección anterior da $\pi_{ij}\le M/n$. Como cada fila tiene masa uno,
$\sum_{ij}\pi_{ij}=n$, y por Cauchy--Schwarz
\[
 1=\frac{(\sum_{ij}\pi_{ij})^2}{n^2}
 \le \sum_{ij}\pi_{ij}^2
 \le \frac{M}{n}\sum_{ij}\pi_{ij}=M .
\]
Por eso la suma de los cuadrados de los pesos permanece en la escala de la constante.
El paso al perfil, (LC3)--(LC4), identifica su límite con $I_2$ y da así el límite $3I_2$ del término local.

Las Secciones~\ref{sec:D01} y~\ref{sec:D02} calculan la contribución de la entropía mediante su
formulación dual. Para el pseudodeterminante hay que identificar tanto el límite del operador
(Sección~\ref{sec:D03}) como el de su traza; para ello se usa el determinante regularizado
$\det_2$ de~\cite[cap.~9]{Simon} (Sección~\ref{sec:D04}). La Sección~\ref{sec:D05} reúne estas
contribuciones y aplica la fórmula de Stirling.
